% ======================================================================
%  sec23.tex —— 第 23 节 有限型态射与有限表现态射
%  底本：原书 pp.111–113（印刷页 = PDF 页）
%  原书标题：23. Morphisms of finite type and finite presentation
% ======================================================================

\section{有限型态射与有限表现态射}\label{sec:23}

% ---------------------------------------------------------------- p.111
\subsection{有限型态射}\label{sec:23.1}

回忆：环同态 $A \to B$ 称为\emph{有限型的}，若 $B$ 在 $A$ 上有限生成，即 $B$ 作为 $A$-代数同构于某个 $A[x_1,\dots,x_n]$ 的商。

\begin{definition}\label{def:23.1.1}
设 $f\colon X \to Y$ 为概形态射。称 $f$ \emph{局部有限型}，若对任意满足 $f(U) \subseteq V$ 的仿射开子概形 $U \subseteq X$ 与 $V \subseteq Y$，诱导的环同态 $\cO_Y(V) \to \cO_X(U)$ 都有限型。

称 $f$ 为\emph{有限型}，若它局部有限型且拟紧。
\end{definition}

\begin{yizhu}
有限型是「局部有限型加拟紧」，而拟紧这一半并非装饰。取可数多个仿射开片的不交并 $\coprod_{n \ge 1} \Spec(k) \to \Spec(k)$：它在每一点的附近都局部有限型，整体却不拟紧，故不有限型。Stacks Project 同样把这两个词分开定义。
[来源：Stacks Project, Tag 01T1（Definition 29.15.1，有限型态射），\url{https://stacks.math.columbia.edu/tag/01T1}（2026-11-17 访问）。]
\end{yizhu}

\begin{lemma}\label{lem:23.1.2}
设 $f\colon X \to Y$ 为概形态射。下列条件等价：
\begin{enumerate}[label=(\roman*),leftmargin=2em,itemsep=0.25ex]
  \item 态射 $f$ 局部有限型；
  \item 对任意点 $x \in X$，存在 $x$ 的仿射开邻域 $U$ 与仿射开子概形 $V \subseteq Y$，使 $f(U) \subseteq V$，且诱导的环同态 $\cO_Y(V) \to \cO_X(U)$ 有限型。
\end{enumerate}
\end{lemma}

该引理来自下述范式。

\begin{exercise}\label{exc:23.1.3}
设 $P$ 是 $\CRing$ 中态射的一个性质，满足：
\begin{enumerate}[label=(\roman*),leftmargin=2em,itemsep=0.25ex]
  \item 对任意 $A \to B$ 与 $a \in A$，有 $P(A \to B) \Rightarrow P(A_a \to B_a)$；
  \item 对任意 $A, B \in \CRing$、$a \in A$、$b \in B$ 与 $A_a \to B$，有 $P(A_a \to B) \Rightarrow P(A \to B_b)$；
  \item 对任意 $A \to B$ 与 $b_1, \dots, b_n \in B$，若 $(b_1, \dots, b_n) = B$ 且对每个 $i$ 都有 $P(A \to B_{b_i})$，则 $P(A \to B)$。
\end{enumerate}
则对概形间的态射 $f\colon X \to Y$，下列条件等价：
\begin{enumerate}[label=(\roman*),leftmargin=2em,itemsep=0.25ex]
  \item 对任意满足 $f(U) \subseteq V$ 的仿射开概形 $U \subseteq X$ 与 $V \subseteq Y$，有 $P(\cO_Y(V) \to \cO_X(U))$；
  \item 对任意点 $x \in X$，存在 $x$ 的仿射开邻域 $U$ 与仿射开子概形 $V \subseteq Y$，使 $f(U) \subseteq V$ 且 $P(\cO_Y(V) \to \cO_X(U))$。
\end{enumerate}
\end{exercise}

\begin{corollary}\label{cor:23.1.4}
态射 $\Spec(B) \to \Spec(A)$ 有限型当且仅当 $A \to B$ 有限型。
\end{corollary}

下列结果由引理 \ref{lem:23.1.2} 容易得到。

\begin{corollary}\label{cor:23.1.5}
（局部）闭浸入是（局部）有限型的。
\end{corollary}

\begin{lemma}\label{lem:23.1.6}
考虑概形的（局部）有限型态射构成的类。
\begin{enumerate}[label=(\roman*),leftmargin=2em,itemsep=0.25ex]
  \item 设 $X \xrightarrow{f} Y \xrightarrow{g} Z$ 为态射链，$f$ 与 $g$ 都（局部）有限型。则 $g \circ f$ 亦然。
  \item 设 $X \xrightarrow{f} Y \xrightarrow{g} Z$ 为态射链，$g \circ f$ 局部有限型。则 $f$ 局部有限型。
  \item[(ii')] 设 $X \xrightarrow{f} Y \xrightarrow{g} Z$ 为态射链，$g \circ f$ 有限型且 $g$ 拟分离。则 $f$ 有限型。
  \item 概形的（局部）有限型态射对基变换稳定。
  \item 「（局部）有限型」在目标上是局部的。
  \item 「局部有限型」在源上是局部的。
\end{enumerate}
\end{lemma}

% ---------------------------------------------------------------- p.112
下列结果来自 Hilbert 基定理。

\begin{lemma}\label{lem:23.1.7}
设 $f\colon X \to Y$ 为概形态射。若 $Y$ （局部）诺特且 $f$ （局部）有限型，则 $X$ （局部）诺特。
\end{lemma}

\begin{yizhu}
Stacks Project 把这条引理的代数内核拆成两条并列陈述：诺特环上的有限生成代数仍诺特；诺特环的局部化仍诺特。前者是 Hilbert 基定理，后者靠理想与局部化的交换性。几何侧的「局部诺特底加局部有限型态射」恰好逐片调用这两条。
[来源：Stacks Project, Tag 00FN（Lemma 10.31.1，诺特性对有限型扩张与局部化稳定），\url{https://stacks.math.columbia.edu/tag/00FN}（2026-11-17 访问）。]
\end{yizhu}

\begin{corollary}\label{cor:23.1.8}
设 $\cC$ 为 $\Sch$ 中由（局部）诺特概形与（局部）有限型态射构成的子范畴。则 $\cC$ 有纤维积。
\end{corollary}

\subsection{有限表现态射}\label{sec:23.2}

回忆：环同态 $A \to B$ 称为\emph{有限表现的}，若 $B$ 作为 $A$-代数同构于某个 $A[x_1,\dots,x_n]$ 模一个有限生成理想所得的商。

\begin{definition}\label{def:23.2.1}
设 $f\colon X \to Y$ 为概形态射。称 $f$ \emph{局部有限表现}，若对任意满足 $f(U) \subseteq V$ 的仿射开子概形 $U \subseteq X$ 与 $V \subseteq Y$，诱导的环同态 $\cO_Y(V) \to \cO_X(U)$ 都有限表现。

称 $f$ 为\emph{有限表现}，若它局部有限表现、拟紧且拟分离。
\end{definition}

\begin{yizhu}
有限表现比有限型多出两件事：定义理想须有限生成（局部条件），整体上还须拟分离。拟分离不可省——开浸入总局部有限表现，但只有拟紧时才有限表现（推论 \ref{cor:23.2.4}）。Stacks Project 的定义同样把这三个条件并列，缺一不可。
[来源：Stacks Project, Tag 01TP（Definition 29.22.1，有限表现态射），\url{https://stacks.math.columbia.edu/tag/01TP}（2026-11-17 访问）。]
\end{yizhu}

我们有引理 \ref{lem:23.1.2} 的如下类比。

\begin{lemma}\label{lem:23.2.2}
设 $f\colon X \to Y$ 为概形态射。下列条件等价：
\begin{enumerate}[label=(\roman*),leftmargin=2em,itemsep=0.25ex]
  \item 态射 $f$ 局部有限表现；
  \item 对任意点 $x \in X$，存在 $x$ 的仿射开邻域 $U$ 与仿射开子概形 $V \subseteq Y$，使 $f(U) \subseteq V$，且诱导的环同态 $\cO_Y(V) \to \cO_X(U)$ 有限表现。
\end{enumerate}
\end{lemma}

\begin{corollary}\label{cor:23.2.3}
态射 $\Spec(B) \to \Spec(A)$ 有限表现当且仅当 $A \to B$ 有限表现。
\end{corollary}

\begin{corollary}\label{cor:23.2.4}
开浸入局部有限表现。拟紧的开浸入有限表现。
\end{corollary}

\begin{warning}\label{warn:23.2.5}
闭浸入未必局部有限表现。
\end{warning}

\begin{yizhu}
反例藏在定义理想里。取 $\pi\colon \Spec(k) \to \Spec(k[x_1,x_2,\dots])$，由 $x_i \mapsto 0$ 给出。该闭浸入的定义理想是极大理想 $\mathfrak m = (x_1,x_2,\dots)$，它不有限生成，故 $\cO(\Spec(k[x_1,x_2,\dots])) \to k$ 不有限表现。闭浸入有限表现，当且仅当定义理想有限生成。
[译者按（上下文推断，无外部依据）：这是「闭浸入未必有限表现」的标准反例。]
\end{yizhu}

我们有引理 \ref{lem:23.1.6} 的如下类比。

\begin{lemma}\label{lem:23.2.6}
考虑概形的（局部）有限表现态射构成的类。
\begin{enumerate}[label=(\roman*),leftmargin=2em,itemsep=0.25ex]
  \item 设 $X \xrightarrow{f} Y \xrightarrow{g} Z$ 为态射链，$f$ 与 $g$ 都（局部）有限表现。则 $g \circ f$ 亦然。
  \item 设 $X \xrightarrow{f} Y \xrightarrow{g} Z$ 为态射链，$g \circ f$ 局部有限表现且 $g$ 局部有限型。则 $f$ 局部有限表现。
  \item[(ii')] 设 $X \xrightarrow{f} Y \xrightarrow{g} Z$ 为态射链，$g \circ f$ 有限表现且 $g$ 局部有限型、拟分离。则 $f$ 有限表现。
  \item 概形的（局部）有限表现态射对基变换稳定。
  \item 「（局部）有限表现」在目标上是局部的。
  \item 「局部有限表现」在源上是局部的。
\end{enumerate}
\end{lemma}

% ---------------------------------------------------------------- p.113
关于（局部）有限型态射与（局部）有限表现态射，我们还有如下显然的结果。

\begin{lemma}\label{lem:23.2.7}
设 $f\colon X \to Y$ 为态射。
\begin{enumerate}[label=(\roman*),leftmargin=2em,itemsep=0.25ex]
  \item 若 $f$ （局部）有限表现，则 $f$ （局部）有限型。
  \item 若 $f$ （局部）有限型且 $Y$ 局部诺特，则 $f$ （局部）有限表现。
\end{enumerate}
\end{lemma}

\begin{yizhu}
Stacks Project 把局部情形与整体情形分列，且不额外假设拟分离。原因是引理 \ref{lem:23.1.7}：底概形局部诺特、态射有限型时 $X$ 局部诺特，再由引理 \ref{lem:22.3.2}，$X \to Y$ 自动拟分离。第 22 节与本节在此接上。
[来源：Stacks Project, Tag 01TX（Lemma 29.22.9，局部诺特底上的有限型即有限表现），\url{https://stacks.math.columbia.edu/tag/01TX}（2026-11-17 访问）。]
\end{yizhu}

\subsection{诺特逼近}\label{sec:23.3}

诺特逼近是一种技巧：只要把注意力限制在\emph{有限表现}的 $S$-概形上，就能把关于一般底概形 $S$ 的问题化归为 $S$ 诺特的情形。这一技巧隐含于 [Gro66, Chapter 8]，并在 [TT13, Appendix C] 中被明确提出。为简单起见，我们只陈述其如下简化版本。

\begin{theorem}\label{thm:23.3.1}
设 $A \simeq \colim_{i \in I} A_i$ 是 $\CRing$ 中的滤过余极限。对任意 $B \in \CRing$，记 $\Sch^{\mathrm{fp}}_B \subseteq \Sch_B$ 为在 $\Spec(B)$ 上有限表现的概形构成的全子范畴。则函子
\[
\Sch^{\mathrm{fp}}_{A_i} \lto \Sch^{\mathrm{fp}}_A, \qquad Y \longmapsto Y \times_{\Spec(A_i)} \Spec(A)
\]
诱导等价
\[
\colim_{i \in I} \Sch^{\mathrm{fp}}_{A_i} \xrightarrow{\ \simeq\ } \Sch^{\mathrm{fp}}_A.
\]
\end{theorem}

\begin{yizhu}
定理可读作「有限表现即紧对象」。Stacks Project 证明：$f$ 局部有限表现，当且仅当对任意滤过系统 $T_i$，有 $\Mor_S(\lim_i T_i, X) = \colim_i \Mor_S(T_i, X)$。诺特逼近正是这条余极限交换性在「底概形也一起变动」时的整体形式，故必须把对象限制在有限表现一侧。
[来源：Stacks Project, Tag 01ZC（Proposition 32.6.1，有限表现与滤过余极限交换），\url{https://stacks.math.columbia.edu/tag/01ZC}（2026-11-17 访问）。]
\end{yizhu}

具体说来，上述定理给出下述结论。

\begin{corollary}\label{cor:23.3.2}
设 $A \simeq \colim_{i \in I} A_i$ 为 $\CRing$ 中的滤过余极限。
\begin{enumerate}[label=(\roman*),leftmargin=2em,itemsep=0.25ex]
  \item 对任意概形 $X \in \Sch^{\mathrm{fp}}_A$，存在指标 $i \in I$ 与 $X_i \in \Sch^{\mathrm{fp}}_{A_i}$，使
  \[
  X \simeq X_i \times_{\Spec(A_i)} \Spec(A).
  \]
  \item 对任意态射 $f\colon X \to X'$（$X, X' \in \Sch^{\mathrm{fp}}_A$），存在指标 $i \in I$ 与态射 $f_i\colon X_i \to X'_i$（$X_i, X'_i \in \Sch^{\mathrm{fp}}_{A_i}$），使 $f$ 可等同于 $f_i$ 沿 $\Spec(A) \to \Spec(A_i)$ 的基变换。
\end{enumerate}
\end{corollary}

\begin{example}\label{ex:23.3.3}
注意任意交换环 $R$ 都可写成有限生成 $\Z$-代数的滤过余极限。实践中，把上述定理用于这样的余极限，即可把关于有限表现 $R$-概形的问题化归为关于有限型 $\Z$-概形的问题。
\end{example}

\begin{yizhu}
Stacks Project 把这一观察推广为：任意拟紧拟分离概形都是有限型 $\Z$-概形的逆向极限。原文只取仿射情形，即 $R$ 为有限生成 $\Z$-代数的滤过余极限；一般情形须把各仿射片沿交叠粘起来，这一步用到拟紧与拟分离，也解释了定理为何只谈有限表现的 $S$-概形。
[来源：Stacks Project, Tag 01ZA（Proposition 32.5.4，绝对诺特逼近），\url{https://stacks.math.columbia.edu/tag/01ZA}（2026-11-17 访问）。]
\end{yizhu}

\begin{example}\label{ex:23.3.4}
设 $S$ 为概形，$s \in S$。注意 $\cO_{S,s} \simeq \colim \cO_S(U)$，其中 $U$ 取遍 $s$ 的仿射开邻域。实践中，把上述定理用于这样的余极限，即可把关于有限表现 $\cO_{S,s}$-概形的构造延拓到 $s$ 的某个开邻域上。
\end{example}
